\subsection{COMPACT SETS IN A UNIFORM SPACE}

The following proposition for an arbitrary uniform space is a sharper form of Proposition 2 of Chapter I, § 9, no. 2 for compact spaces.

PROPOSITION 4. In a uniform space X, let A be a compact set and B a closed set such that A∩B=∅. Then there is an entourage V of X such that V(A) and V(B) do not intersect.

If the proposition were false, then none of the sets \(A \cap \mathring{V}(B)\) , where V runs through the set of symmetric entourages of X, would be empty; hence these sets would form a filter base on A, which would have a cluster point \(x_0 \in A\) . Hence for each symmetric entourage \(V\) of \(X\) , \(\dot{V}(x_0)\) would meet \(B\) and therefore, as \(B\) is closed, we should have \(x_0 \in B\) , contrary to hypothesis.

COROLLARY. Let A be a compact set in a uniform space X; then as V runs through the set of entourages of X, the sets V(A) form a fundamental system of neighbourhoods of A.

Let U be any open neighbourhood of A, then B = CU is closed and does not meet A; hence, by Proposition 4, there is an entourage V such that \(\mathrm{V}(\mathrm{A}) \cap \mathrm{V}(\mathrm{B}) = \varnothing\) , and therefore \(\mathrm{V}(\mathrm{A}) \subset \mathrm{U}\) .

